\documentclass[12pt]{article}
\usepackage{amsmath,amssymb}

\begin{document}
\section*{{2023 AIME II Problems/Problem 7}}
Source: AoPS Wiki

\subsection*{Solution}
Note that the condition is equivalent to stating that there are no 2 pairs of oppositely spaced vertices with the same color. Case 1: There are no pairs. This yields $2$ options for each vertices 1-6, and the remaining vertices 7-12 are set, yielding $2^6=64$ cases. Case 2: There is one pair. Again start with 2 options for each vertex in 1-6, but now multiply by 6 since there are 6 possibilities for which pair can have the same color assigned instead of the opposite. Thus, the cases are: $2^6*6=384$ case 3: There are two pairs, but oppositely colored. Start with $2^6$ for assigning 1-6, then multiply by 6C2=15 for assigning which have repeated colors. Divide by 2 due to half the cases having the same colored opposites. $2^6*15/2=480$ It is apparent that no other cases exist, as more pairs would force there to be 2 pairs of same colored oppositely spaced vertices with the same color. Thus, the answer is: $64+384+480=\boxed{928}$ ~SAHANWIJETUNGA

\end{document}
